%! vim: spl=en:
\documentclass[a0paper,portrait,fontscale=.25]{baposter}

\usepackage{relsize}
\usepackage[utf8]{inputenc}
\usepackage{PTSans}
\renewcommand{\familydefault}{\sfdefault}

\begin{document}

\background{%
%\includegraphics{paticky_tisk_up_01.pdf}
}

\def\bp{\parindent0pt\parskip=\smallskipamount\raggedright}
\definecolor{main}{cmyk}{0.06,0.26,0.91,0}
\definecolor{other}{cmyk}{0.94,0.74,0.09,0}
\definecolor{muni}{RGB}{0,26,211}
\definecolor{vscht}{RGB}{205,84,55}

\def\alert#1{\textbf{\color{vscht}#1}}
\def\qrcite#1{\lower.9ex\hbox{{\includegraphics[height=2.8ex]{#1}}}}

\begin{poster}{
	grid=false,
	columns=6,
	background=user,
	borderColor=gray,
	headerFontColor=white,
	headerColorOne=vscht,
	headerColorTwo=muni,
	textborder=rectangle,
%	headershape=roundedright,
%	headerborder=none,
%	boxColorOne=main,
%	boxColorTwo=main,
	headerfont=\Large\bfseries,
	headerborder=open,
	boxshade=none,
}
{ % Eyecatcher
\includegraphics[height=25mm]{poster-qr}
}
{%
\huge
The Chicken-and-Egg Problem of Landmark-Driven Molecular Dynamics: Are Random Landmarks Useful?
}
{
\large\vrule height2ex width0pt {\color{muni}Aleš Křenek, Jana Hozzová, Jaroslav Oľha, Martin Kurečka,}
{\color{vscht} Dalibor Trapl, Vojtěch Spiwok}
}
{ % logo
%\hfill\includegraphics[height=10mm]{mubeamer-ics-english-color}
%\hfill\includegraphics[height=10mm]{logoUCT_basic}
}

\headerbox{Metadynamics with PCV}{name=pcv,column=0,row=0,span=2}{
\bp
Molecular dynamics of proteins can be guided to \alert{explore wider range} of conformational space (folding paths in particular) with \alert{biased potential} built on \alert{path collective variables}.

PCVs are derived from a \alert{set of landmark structures} which approximate the desired trajectory
%\penalty-10000
[ \qrcite{from-a-to-b} \qrcite{meta-isomap} ].

The technique can \alert{reduce the time} of the MD simulation dramatically. On the other hand, the choice of landmarks is the core of the \alert{\emph{chicken and egg} problem} – if we know the landmarks along the trajectory, what would be the reason of recomputing the same trajectory?
% (Well, things are not so simple, but this is the core.)
}

\headerbox{}{name=bigpic,column=0,span=6,below=pcv,boxheaderheight=0pt,textborder=none,headerborder=none}
{
\includegraphics[width=\hsize]{bigpic2}
}

\headerbox{Randomly Generated Landmarks}{name=vystupy,column=2,span=4,above=bigpic}{
\bp
We aim at \alert{computing the trajectories de novo}, without their prior knowledge. The essential steps are:
\begin{enumerate}
\item Generate a~set of several hundreds to thousands of barely feasible landmarks 
by \alert{random twisting of peptide bonds} in the subject protein
\item \alert{Minimize by steepest descent} in vacuo, using simple force field (Amber99) to resolve unrealistic properties
\item Discard structures with \alert{too high energy} or failed minimization
\item Compute low-dimensional embedding of 3N atomic coordinates to define \alert{collective variables}
\item Train feed-forward \alert{neural network} to estimate CVs from coordinates [ \qrcite{anncolvar} ]
\item Use the network outputs to \alert{generate bias potential} in metadynamics simulation
\end{enumerate}
}


\headerbox{Implementation}{name=impl,column=0,below=bigpic,span=4}{
\bp
The whole workflow is implemented in Python as a~\alert{Jupyter notebook} using few standard biochemical 
packages.

The neural network training uses our \alert{Anncolvar} package on top of Keras.

Core of the MD simulation is \alert{Gromacs}.

Estimation of the CVs by neural network is done by our \alert{optimized module in Plumed}.
}

\headerbox{}{name=logos,column=0,span=6,boxheaderheight=0pt,below=impl,textborder=none,headerborder=none}{
\bp
\includegraphics[height=10mm]{logoUCT_basic}
\hfill\includegraphics[height=10mm]{mubeamer-ics-english-color}
}

\headerbox{Acknowledgement}{name=ack,column=4,below=bigpic,span=2,above=logos}{
\bp
\footnotesize
\raggedright
The work was supported by the Grant Agency of the Czech Republic, project no.\ GA19-16857S.

Computational resources were provided by the project LM2018140 supported by the Ministry of \hbox{Education}, Youth and Sports of the Czech Republic.
}


\end{poster}
\end{document}
